% SPDX-License-Identifier: CC-BY-4.0
% Copyright (c) 2026 Martin Schröder <info@swedishembedded.com>

\section{Training a 27B model on one card}\label{sec:substrate}

Training a 27-billion-parameter model on one 95~GiB card needs the frozen base in a
reduced-precision tier, activation recomputation, a loss that never materialises the
full logits, an attention kernel that supports the model's head dimension, and a loader
that never holds the checkpoint in host memory. This section states what Brain's main
tree contains for each, what its qualification runs measured for bf16 LoRA and NF4
QLoRA, and what a later throughput measurement of the same architecture shows. It makes
no claim that a trained 27B adapter changes what the model generates; that is the
purpose of experiment CE6 (\cref{sec:pending}).

\subsection{What the tree contains}

The Splinter trainer's seam is the Qwen3-family trainer.\cl{B01} In that trainer the
frozen base can be held in bf16,\cl{B02} each layer is recomputed in the backward pass
from its saved residual,\cl{B03} and the logits are sized in row chunks with a chunked
weighted cross-entropy.\cl{B04} Brain has a causal grouped-query flash-attention forward
and backward supporting head dimensions up to 256, with finite-difference gradient
checks at head dimensions 40, 128 and 256.\cl{B05,B10} For the 27B model's crate it has
a streaming FP8 source that dequantises in 128-row bands and a base-dtype enumeration
with F32, bf16 and NF4 members,\cl{B06} and an NF4 weight type with a 16-entry lookup
table.\cl{B07} The numerical gates are \code{make test/gradcheck} (finite-difference
checks of the Qwen3.5 architecture including LoRA) and \code{make test/parity}, which
runs the checks on the CPU backend.\cl{B08,B09}

\subsection{Qualification runs}\label{sec:qual}

Brain's qualification at commit \mEFourCommit{} ran the 27B architecture on one 95~GiB
card; \cref{tab:qual} lists the runs. Probes use a randomly initialised model of the real
width and depth (64 layers), recomputation, rank-16 LoRA and one step, so their loss is not
a quality number. With an NF4 base, 8192 tokens fit in
\mEFourProbeNfFourEightOneNineTwoPeakMib{}~MiB at
\mEFourProbeNfFourEightOneNineTwoTokensPerS{} tokens/s;\cl{M11} with a bf16 base,
8192 tokens took \mEFourProbeBfOneSixEightOneNineTwoPeakMib{}~MiB and 16384 tokens
\mEFourProbeBfOneSixOneSixThreeEightFourPeakMib{}~MiB, both at
\mEFourProbeBfOneSixOneSixThreeEightFourTokensPerS{} tokens/s.\cl{M12,M13}
Through the command line the real FP8 checkpoint trained one step with an NF4 base and
with a bf16 base with a host peak of \mEFourCliNfFourOnestepHostRssMib{} and
\mEFourCliBfOneSixOnestepHostRssMib{}~MiB respectively\cl{M14,M15} and, with an NF4 base,
40 steps drove the loss on eight training examples from
\mEFourCliNfFourFourZerostepLossFirst{} to \mEFourCliNfFourFourZerostepLossLast{}.\cl{M16}
That is a statement about fitting the training examples, not about generalisation: the
validation examples were two of the training examples. The command line's default
batch ran the device out of memory at 1024-token rows.\cl{M17}

\begin{table}[t]
\centering
\small
\input{generated/e4_table}
\caption{Qualification runs of the 27B architecture on one card (Brain's qualification
logs, copied unmodified apart from path names; their values, digests and commands are in \code{data/measurements.json}).}\label{tab:qual}
\end{table}

\subsection{Throughput as measured on a later Brain revision}\label{sec:throughput}

The qualification rates above are those of the commit it ran at. A later probe of the
same architecture (random weights, 64 layers, rank 16, NF4 base, 8192 tokens, recomputation)
on the idle card measured \mThrNfFourEightOneNineTwoTokensPerS{} tokens/s at steps two and
three, a step of \mThrNfFourEightOneNineTwoStepS{}~s and a peak of
\mThrNfFourEightOneNineTwoPeakMib{}~MiB.\cl{M38} Brain's own record of the real checkpoint
at 8192 tokens states 635 tokens/s with an NF4 base, and 640 tokens/s with a bf16 base at
79.8~GiB;\cl{B13,B14} we did not rerun those two, and the same record states that
1000 tokens/s was not reached.\cl{B15} At the measured probe rate one million trained tokens
take \mThrHoursPerMillionTokens{}~h of card time, computed as $10^6/(622\cdot3600)$.\cl{C12}
We state the measurement and not a target; a rate is only valid for the revision, the
card and the sequence length it was taken at.

\subsection{What this does not show}

That a trained 27B adapter changes what the model generates. The qualification logs predate
a generation path for such an adapter; Brain has since added one (the \code{sample-adapter}
verb of its \code{qwen35} command)\cl{B12} and keeps its own record of a run, which we did
not rerun or verify here. No run beyond 16384 tokens and no NF4 run at 16384 tokens on the
qualification commit exists.
